%% ================================================================
%% DSTP Non-Reducibility Remark (final version)
%% To be inserted in paper1.tex immediately after
%% \begin{remark}[DSTP heuristics] and before
%% \begin{theorem}[Conditional frame stability].
%% ================================================================

\begin{remark}[Non-reducibility of \textup{DSTP} --- constructive witness]
\label{rem:ax5_independence}
DSTP (Definition~\ref{def:dstp}) is the concrete realisation of the
abstract interface hypothesis \textup{AX5}: while \textup{AX1--AX4}
control the \emph{analytic} structure of the system (symbol class,
spectral gaps, quadrature order, uniform PSWF bounds), DSTP controls
the \emph{geometric} compatibility of the discretisation.
The two layers are logically independent within this framework,
as the following constructive witness shows.

\smallskip
\noindent\textbf{Witness rule.}
Let $\{p_j = -T + 2Tj/N\}_{j=1}^N$ be the equidistant Riemann rule
with weights $w_j = 2T/N$ and mesh $h = 2T/N$.
This rule satisfies:
\begin{enumerate}[(i)]
\item \textup{AX3} (first-order quadrature consistency,
  Assumption~\ref{ass:quad_first_order}), with $C_Q = 1$;
\item \textup{AX4} (uniform PSWF bounds \cite{BonamiKaroui2014}),
  since $\|\psi_n^{(c)}\|_{L^\infty}$ and
  $\|(\psi_n^{(c)})'\|_{L^\infty}$ are intrinsic to the PSWFs and
  independent of the node locations;
\item \textup{AX1--AX2} (symbol class and spectral gap), which are
  consequences of \cite{PaperIV,OsipovRokhlinXiao} alone.
\end{enumerate}

\smallskip
\noindent\textbf{Spectral concentration of $f_{mn}$.}
Since $\hat{f}_{mn} = \hat{\psi}_m^{(c)} \ast \hat{\psi}_n^{(c)}$
(convolution theorem), and each factor $\hat{\psi}_k^{(c)}$ is
exponentially concentrated on $[-\omega/\pi, \omega/\pi]$
\cite{BonamiKaroui}, the convolution is exponentially concentrated
on $[-2\omega/\pi, 2\omega/\pi] = [-4c/T, 4c/T]$:
\[
  \int_{|\xi| > 4c/T} |\hat{f}_{mn}(\xi)|^2\,d\xi \le C e^{-\alpha c},
\]
with $\alpha, C > 0$ independent of $m, n, c$.
In particular, $\hat{f}_{mn}(k/h) = O(e^{-\alpha c})$ for $|k| \ge 2$,
since $2/h \sim 2N/(2T) \ge 2c/T \cdot (2/\pi) > 4c/T$ for
$N = \lfloor 2c/\pi\rfloor$.
Hence the aliasing sum is dominated by the $k = 1$ term.

\smallskip
\noindent\textbf{DSTP fails.}
By the Poisson summation formula,
\[
  h\sum_{j=1}^N f_{mn}(p_j)
  = \int_{-T}^T f_{mn}\,dx
  + \sum_{k \ne 0} \hat{f}_{mn}(k/h).
\]
Since $1/h \sim c/(2T)$ lies strictly within the spectral
concentration band $[0, 4c/T]$ of $\hat{f}_{mn}$, the $k=1$
aliasing term is not exponentially small.
Accumulating over all off-diagonal pairs,
\[
  \sum_{m \ne n} |E_{mn}^{\mathrm{Riemann}}|^2
  \;\ge\; \sum_{m \ne n} |\hat{f}_{mn}(1/h)|^2 - O(N^2 e^{-\alpha c})
  \;\ge\; c_0\,(hc)^2 N,
\]
with $c_0 > 0$ independent of $c \ge c_0$, from the non-negligible
aliasing mass at frequency $1/h$. The off-diagonal mass is not
concentrated on the diagonal: for $N \sim c$, the aliasing frequency
$1/h$ is active across all index pairs $(m,n)$ with
$m \ne n \in \mathcal{B}$, since the spectral support of $\hat{f}_{mn}$
is $m$- and $n$-independent at leading order.
Therefore:
\[
  \|E_{N,c}^{\mathrm{Riemann}}\|_2
  \;\ge\; \frac{\|E_{N,c}^{\mathrm{Riemann}}\|_F}{\sqrt{N}}
  \;\ge\; c_0^{1/2}\,hc,
\]
with no decay in $|m-n|$, so DSTP fails for all $p > 0$.

\smallskip
\noindent\textbf{Structural interpretation.}
The failure of DSTP for the Riemann rule is not a consequence of
low approximation order or insufficient regularity.
\textup{AX1--AX4} together supply $c$-uniform amplitude control,
spectral gaps, quadrature consistency, and PSWF bounds --- yet none
of these encodes geometric compatibility between the node set
$\{p_j\}$ and the oscillatory eigenstructure of $V_{N,c}$.
DSTP is not an estimate of higher order; it is a condition of a
different type entirely:
\begin{quotation}
  \noindent\textit{DSTP is not an estimate. It is a spectral
  incompatibility condition: the aliasing frequency of the sampling
  rule falls within the spectral concentration band of the PSWF
  product kernel.}
\end{quotation}
This places DSTP in the same logical category as spectral exactness
conditions in Gauss quadrature, sampling compatibility in the
Shannon--Nyquist theorem, and frame reconstruction axioms in
non-uniform sampling theory --- none of which follow from regularity
bounds alone.

\smallskip
\noindent\textbf{Minimality of the axiom system.}
DSTP is not ``stronger'' than \textup{AX1--AX4}; it is
\emph{orthogonal} to them, in the sense of logical independence
of admissible discretisation classes within this framework,
not linear-algebraic orthogonality:
\textup{AX1--AX4} admit any sampling geometry consistent with
regularity and approximation order; DSTP selects those geometries
compatible with the spectral structure of $V_{N,c}$.
No combination of analytic bounds implies geometric compatibility,
and the Riemann-rule witness shows that the converse also fails:
a rule may satisfy all of \textup{AX1--AX4} and fail DSTP.
The axiom system $\{\textup{AX1},\dots,\textup{AX4},\textup{DSTP}\}$
is therefore minimal in the sense of logical independence:
no axiom is derivable from the remaining ones within this framework.

Verification of DSTP for PSWF-adapted (XRY/Gauss--PSWF) quadrature
schemes is the content of~\cite{PaperIIquad}.
\end{remark}
